\documentclass{article}

\usepackage{tikz} 
\usetikzlibrary{automata, positioning, arrows} 

\usepackage{amsthm}
\usepackage{amsfonts}
\usepackage{amsmath}
\usepackage{amssymb}
\usepackage{fullpage}
\usepackage{color}
\usepackage{parskip}
\usepackage{hyperref}
  \hypersetup{
    colorlinks = true,
    urlcolor = blue,       % color of external links using \href
    linkcolor= blue,       % color of internal links 
    citecolor= blue,       % color of links to bibliography
    filecolor= blue,       % color of file links
    }
    
\usepackage{listings}
\usepackage[utf8]{inputenc}                                                    
\usepackage[T1]{fontenc}                                                       

\definecolor{dkgreen}{rgb}{0,0.6,0}
\definecolor{gray}{rgb}{0.5,0.5,0.5}
\definecolor{mauve}{rgb}{0.58,0,0.82}

\lstset{frame=tb,
  language=haskell,
  aboveskip=3mm,
  belowskip=3mm,
  showstringspaces=false,
  columns=flexible,
  basicstyle={\small\ttfamily},
  numbers=none,
  numberstyle=\tiny\color{gray},
  keywordstyle=\color{blue},
  commentstyle=\color{dkgreen},
  stringstyle=\color{mauve},
  breaklines=true,
  breakatwhitespace=true,
  tabsize=3
}

\newtheoremstyle{theorem}
  {\topsep}   % ABOVESPACE
  {\topsep}   % BELOWSPACE
  {\itshape\/}  % BODYFONT
  {0pt}       % INDENT (empty value is the same as 0pt)
  {\bfseries} % HEADFONT
  {.}         % HEADPUNCT
  {5pt plus 1pt minus 1pt} % HEADSPACE
  {}          % CUSTOM-HEAD-SPEC
\theoremstyle{theorem} 
   \newtheorem{theorem}{Theorem}[section]
   \newtheorem{corollary}[theorem]{Corollary}
   \newtheorem{lemma}[theorem]{Lemma}
   \newtheorem{proposition}[theorem]{Proposition}
\theoremstyle{definition}
   \newtheorem{definition}[theorem]{Definition}
   \newtheorem{example}[theorem]{Example}
\theoremstyle{remark}    
  \newtheorem{remark}[theorem]{Remark}

\title{CPSC-354 Report}
\author{Nathan Garcia  \\ Chapman University}

\date{\today} 

\begin{document}

\maketitle

\begin{abstract}
\end{abstract}

\setcounter{tocdepth}{3}
\tableofcontents

\section{Introduction}\label{intro}

\section{Week by Week}\label{homework}

\subsection{Week 1}
\subsubsection{MU Puzzle}

The MU puzzle comes from the book \textit{Gödel, Escher, Bach}. You start with the string \textbf{MI} and the goal is to turn it into \textbf{MU} by following four rules:

\begin{enumerate}
    \item If a string ends with \textbf{I}, you can add a \textbf{U} at the end.  
    \item If a string starts with \textbf{M}, you can copy everything after the M.  
    \item If you see \textbf{III}, you can change it to \textbf{U}.  
    \item If you see \textbf{UU}, you can delete it.  
\end{enumerate}

The puzzle is about seeing if you can reach \textbf{MU} by only using these rules. It is not really about the letters themselves, but about how rules control what strings you can or cannot make.
\subsubsection{HW1}

The MU puzzle comes from the book \textit{Gödel, Escher, Bach}. You start with the string \textbf{MI} and the goal is to turn it into \textbf{MU} by following four rules:

\begin{enumerate}
    \item If a string ends with \textbf{I}, you can add a \textbf{U} at the end.  
    \item If a string starts with \textbf{M}, you can copy everything after the M.  
    \item If you see \textbf{III}, you can change it to \textbf{U}.  
    \item If you see \textbf{UU}, you can delete it.  
\end{enumerate}

At first, I tried small derivations. For example:
\[
\texttt{MI} \;\Rightarrow\; \texttt{MIU} \;\Rightarrow\; \texttt{MIUIU}
\]
or duplicating I’s:
\[
\texttt{MI} \;\Rightarrow\; \texttt{MII} \;\Rightarrow\; \texttt{MIIII}
\]
From \texttt{MIIII}, I can replace \texttt{III} with \texttt{U}, giving \texttt{MUI}, but not \texttt{MU}. Every time, an extra \texttt{I} is left over, and there is no rule that deletes a single \texttt{I}.

\paragraph{Invariant Argument.}  
Let $\#I(w)$ denote the number of \texttt{I}'s in string $w$. If we track $\#I(w) \pmod 3$, we find:
\begin{itemize}
    \item Rule 1: $\#I$ unchanged.  
    \item Rule 2: $\#I$ doubles. Over $\mathbb{Z}/3\mathbb{Z}$, $1 \mapsto 2$, $2 \mapsto 1$, never $0$.  
    \item Rule 3: Removes $3$ I’s, leaving the remainder mod 3 unchanged.  
    \item Rule 4: Deletes U’s only, so $\#I$ unchanged.  
\end{itemize}

We start with \texttt{MI}, which has $\#I=1$. This is congruent to $1 \pmod 3$. Because no rule ever makes $\#I \equiv 0 \pmod 3$, it is impossible to reach a string with $\#I=0$.

\paragraph{Conclusion.}  
The target \texttt{MU} has $\#I=0$, which is divisible by $3$. Since that is unreachable from \texttt{MI}, the puzzle is unsolvable. As a student, the cool part here is that the solution isn’t about brute-force trying rules—it’s about spotting a hidden invariant (the number of I’s mod 3) that blocks the path completely.



\section{Essay}

\section{Evidence of Participation}

\section{Conclusion}\label{conclusion}

\begin{thebibliography}{99}
\bibitem[BLA]{bla} Author, \href{https://en.wikipedia.org/wiki/LaTeX}{Title}, Publisher, Year.
\end{thebibliography}

\end{document}
